\section*{Chapter \ref{ch:g9:metals-acids-corrosion} --- Metals: Reactions with Acids and Corrosion}

\begin{solution}{exo:g9:metals-acids-corrosion:1}
Hydrogen: a lighted splint at the mouth of the tube gives a squeaky pop.
\end{solution}

\begin{solution}{exo:g9:metals-acids-corrosion:2}
Iron, zinc and aluminium react; copper and gold do not.
\end{solution}

\begin{solution}{exo:g9:metals-acids-corrosion:3}
Both oxygen and water.
\end{solution}

\begin{solution}{exo:g9:metals-acids-corrosion:4}
An \omterm{def:g9:ions:ion}{ion} present during a reaction that takes no part in it. With
hydrochloric acid: the chloride \omterm{def:g9:ions:ion}{ion} \ce{Cl-}.
\end{solution}

\begin{solution}{exo:g9:metals-acids-corrosion:5}
\ce{Mg(s) + 2H+(aq) -> Mg^{2+}(aq) + H2(g)}; charge $+2$ on each side.
\end{solution}

\begin{solution}{exo:g9:metals-acids-corrosion:6}
Add sodium hydroxide: a green \omterm{def:g9:ions:precipitate}{precipitate} of \ce{Fe(OH)2} shows
iron(II) \omterm{def:g9:ions:ion}{ions}.
\end{solution}

\begin{solution}{exo:g9:metals-acids-corrosion:7}
Boiling removes the \omterm{def:g4:air-a-mixture-of-gases:air}{air} dissolved in the water, and the oil stops \omterm{def:g4:air-a-mixture-of-gases:air}{air}
from dissolving again: the nail has water but no oxygen.
\end{solution}

\begin{solution}{exo:g9:metals-acids-corrosion:8}
Salt dissolved in the water splashed on the car makes rusting faster.
\end{solution}

\begin{solution}{exo:g9:metals-acids-corrosion:9}
The oxide layer formed sticks to the metal and seals it from the \omterm{def:g4:air-a-mixture-of-gases:air}{air}: the
attack stops at the surface.
\end{solution}

\begin{solution}{exo:g9:metals-acids-corrosion:10}
Paint it, oil or grease it, galvanise it (or make it of stainless steel).
\end{solution}

\begin{solution}{exo:g9:metals-acids-corrosion:11}
\ce{2Al(s) + 6H+(aq) -> 2Al^{3+}(aq) + 3H2(g)}: 2 Al and 6 H on each side,
charge $+6$ on each side. For 200 Al \omterm{def:g7:atoms-and-molecules:atom}{atoms}: $200 \times \frac{3}{2} = 300$
hydrogen \omterm{def:g7:atoms-and-molecules:molecule}{molecules}.
\end{solution}

\begin{solution}{exo:g9:metals-acids-corrosion:12}
The zinc around the scratch is corroded in place of the steel: as long
as zinc remains nearby, the steel is spared.
\end{solution}

\begin{solution}{pb:g9:metals-acids-corrosion:1}
\textbf{1.} \ce{Zn(s) + 2H+(aq) -> Zn^{2+}(aq) + H2(g)}.

\textbf{2.} Hydrogen: a squeaky pop with a lighted splint.

\textbf{3.} Sodium hydroxide gives a white \omterm{def:g9:ions:precipitate}{precipitate} of
\ce{Zn(OH)2}.

\textbf{4.} The chloride \omterm{def:g9:ions:ion}{ions} \ce{Cl-}.

\textbf{5.} $125.6 - 113.5 = \qty{12.1}{g}$.

\textbf{6.} Hydrogen is given off as long as zinc reacts with the acid;
when no zinc is left, the bubbling stops.

\textbf{7.} The iron of the steel would react too, giving hydrogen and
\ce{Fe^{2+}} \omterm{def:g9:ions:ion}{ions}; sodium hydroxide would then give a green \omterm{def:g9:ions:precipitate}{precipitate}.

\textbf{8.} The zinc is attacked in place of the steel and keeps water
and \omterm{def:g4:air-a-mixture-of-gases:air}{air} away from it.

\textbf{9.} $2 \times 10.0 \times 10.0 = \qty{200}{cm^2}$.

\textbf{10.} $12.1 \div 7.13 \approx \qty{1.70}{cm^3}$.

\textbf{11.} $1.70 \div 200 = \qty{0.0085}{cm}$.

\textbf{12.} $0.0085 \times 10\,000 = \qty{85}{\micro m}$: the zinc
layer was about \qty{85}{\micro m} thick.
\end{solution}
